% !TeX spellcheck = ru_RU
% !TEX root = vkr.tex

\section{Заключение} \label{sec:conclusion}

В рамках настоящей работы проведено экспериментальное исследование влияния дообучения большой языковой модели YandexGPT~5.1~Pro методом LoRA~\cite{hu2022lora} на качество классификации инструментов в AI-агенте, функционирующем в домене недвижимости. Исследование выполнено на двух независимых бенчмарках общим объёмом 750~примеров (500 смешанных и 250 реальных) и охватывает 12~моделей с вариацией типа и объёма обучающих данных.

Основные задачи, поставленные в начале исследования, были решены в полном объёме:
\begin{enumerate}
    \item Формализована задача классификации инструментов как задача бинарной классификации с классами GEO (\texttt{ResolveEntities}) и TAG (\texttt{SearchTags}); сформированы два независимых бенчмарка для раздельной и совместной оценки моделей.
    \item Разработан модульный программный пайплайн, обеспечивающий полный цикл~--- от подготовки обучающих данных и запуска дообучения через API YandexGPT до автоматизированного inference и вычисления метрик с визуализацией результатов.
    \item Проведена серия из 12~экспериментов по дообучению с систематической вариацией источника данных (синтетические, реальные, смешанные), классового баланса и объёма выборки (от 50 до ${\sim}\,1500$~примеров).
    \item Выполнен строгий статистический анализ результатов с применением критерия Кохрена~$Q$ для проверки глобальной гетерогенности и попарного критерия Макнемара с поправкой Бонферрони для множественных сравнений.
\end{enumerate}

Главные научные результаты работы заключаются в следующем:
\begin{itemize}
    \item \textbf{Подтверждена гипотеза о стилистической природе выбора инструмента.} Дообучение методом LoRA повышает accuracy с $81{,}8\%$ (базовая модель) до $99{,}4\%$ (лучшая дообученная модель \texttt{train\_mixed\_500}), а значение Macro~$F_1$~--- с $0{,}526$ до $0{,}996$. Прирост accuracy составляет $+17{,}6$~п.\,п., прирост Macro~$F_1$~--- $+0{,}470$. Различия статистически значимы: критерий Кохрена даёт $Q = 348{,}0$, $p < 10^{-6}$, а попарные сравнения по критерию Макнемара подтверждают значимость для всех пар <<базовая~--- дообученная>>. Данный результат согласуется с гипотезой поверхностного выравнивания~\cite{zhou2023lima}: небольшой объём примеров корректного формата оказывается достаточным для радикального улучшения качества, поскольку модель уже обладает необходимыми знаниями о предметной области.

    \item \textbf{Установлена оптимальная стратегия формирования обучающей выборки.} Модель, дообученная на смешанном датасете из 250~реальных и 250~синтетических примеров (\texttt{train\_mixed\_500}), достигает accuracy $99{,}4\%$, превосходя модели, обученные на чисто синтетических данных: $97{,}8\%$ для \texttt{synth\_500} и $95{,}6\%$ для \texttt{synth\_1000}. Увеличение объёма синтетических данных сверх~500 примеров не только не улучшает, но ухудшает результат~--- наблюдается деградация до $95{,}6\%$ при 908~синтетических примерах, что свидетельствует о переобучении на артефакты генерации. Наиболее робастной оказывается именно модель на смешанных данных: её отклонение на реальном бенчмарке составляет лишь $-0{,}6$~п.\,п., тогда как для моделей на чистой синтетике~--- от $-2{,}2$ до $-5{,}4$~п.\,п. Таким образом, качество обучающих данных оказывается важнее их объёма.

    \item \textbf{Обнаружена нестабильность процедуры дообучения.} Три версии модели (v1, v2, v3), дообученные на идентичных данных, демонстрируют разброс accuracy в $5{,}6$~п.\,п. при неконтролируемой стохастичности, обусловленной отсутствием фиксации random seed на стороне API. Этот результат указывает на необходимость учёта случайной вариативности при оценке моделей и ставит под вопрос надёжность одиночных прогонов дообучения.

    \item \textbf{Выявлен эффект консервации ошибок и влияние safety-фильтра.} $73{,}6\%$~примеров бенчмарка классифицируются одинаково всеми моделями, что указывает на устойчивое ядро <<лёгких>> и <<трудных>> запросов. Два запроса, содержащих упоминание Крыма, систематически отклоняются safety-фильтром YandexGPT вне зависимости от дообучения, что демонстрирует ограничения адаптации моделей с жёсткой модерацией контента.
\end{itemize}

Практическая значимость работы определяется следующим. Во-первых, разработанный программный пайплайн (описанный в разделе~\ref{sec:implementation} и предшествующей работе~\cite{goroshko2025otchet}) может быть использован для дообучения и систематической оценки моделей YandexGPT в произвольных задачах классификации. Во-вторых, сформулированы конкретные рекомендации по формированию обучающих данных: использование смешанных выборок (реальные~+ синтетические данные) предпочтительнее наращивания объёма чисто синтетических примеров, что позволяет снизить затраты на ручную разметку без потери качества. В-третьих, показано, что даже для коммерческих моделей, доступных только через API, метод LoRA обеспечивает переход от неприемлемого уровня качества ($81{,}8\%$) к практически безошибочной классификации ($99{,}4\%$).

\subsection{Направления дальнейших исследований}

Проведённое исследование открывает ряд перспективных направлений для продолжения работы:
\begin{itemize}
    \item \textbf{Переход к multi-label классификации.} В текущей постановке каждый запрос относится ровно к одному классу. Однако на практике встречаются запросы, содержащие одновременно географическую и атрибутивную составляющие (например, <<трёхкомнатная квартира на Васильевском острове>>). Расширение модели до multi-label классификации позволит обрабатывать такие случаи корректно.

    \item \textbf{Расширение бенчмарка и ревалидация разметки.} Текущие бенчмарки общим объёмом 750~примеров (500~в смешанном и 250~в реальном), хотя и достаточны для статистического анализа, могут быть расширены для повышения мощности критериев и выявления более тонких различий между моделями. Кроме того, целесообразна независимая экспертная ревалидация разметки для устранения возможных ошибок в эталонных метках.

    \item \textbf{Стабилизация процедуры дообучения.} Обнаруженный разброс в $5{,}6$~п.\,п. между версиями модели мотивирует исследование методов стабилизации: фиксацию random seed (при наличии такой возможности в API), ансамблирование нескольких прогонов дообучения, а также применение техник регуляризации.

    \item \textbf{Исследование альтернативных PEFT-методов.} В данной работе использован исключительно метод LoRA. Сравнение с другими подходами к параметрически эффективному дообучению~--- QLoRA, Prompt Tuning~\cite{lester2021power}, Adapter Tuning~--- позволит определить, является ли LoRA оптимальным выбором для задачи tool selection.

    \item \textbf{Обобщение на другие LLM.} Все эксперименты проведены на модели YandexGPT~5.1~Pro. Проверка воспроизводимости результатов на моделях иных архитектур и провайдеров позволит установить, в какой мере полученные закономерности являются универсальными, а в какой~--- специфичными для конкретной модели.
\end{itemize}

Таким образом, настоящая работа вносит вклад в понимание механизмов адаптации больших языковых моделей к задаче выбора инструментов и демонстрирует, что дообучение методом LoRA на относительно небольшой, но качественно сформированной обучающей выборке является эффективной стратегией повышения качества классификации в AI-агентах.
